% \iffalse meta-comment
% LTeX: language=de-DE
%
%  TUDCD-Script -- Corporate Design of Technische Universität Dresden
% ----------------------------------------------------------------------------
%
%  Copyright (C) TU Dresden <tudlatex@tu-dresden.de>, 2025
%
% ----------------------------------------------------------------------------
%
%  This work may be distributed and/or modified under the conditions of the
%  LaTeX Project Public License, either version 1.3c of this license or
%  any later version. The latest version of this license is in
%    http://www.latex-project.org/lppl.txt
%  and version 1.3c or later is part of all distributions of
%  LaTeX version 2008-05-04 or later.
%
%  This work has the LPPL maintenance status "maintained".
%
%  The current maintainer and author of this work is Yoshi Diepelt.
%
% ----------------------------------------------------------------------------
% \fi
% \selectlanguage{ngerman}
%
% \subsection{\texttt{test-compile}: Übersetzungstest für Fehlerhafte und korrekte Übersetzungen}
%
% Diese Folge an Test ist genau dann erfolgreich, wenn die Kompilation
% ohne Fehler abbricht und eine PDF generiert.
%
% Mit dem folgenden Guard wird die bereitgestellte Testsuite von \dpkg{l3build}
% in alle Tests dieser Suite eingebunden.
% Da die Batchdatei zum erstellen des Tests zerstückelt geschrieben wird,
% müssen Guards um die Präambel der Tests geschrieben werden.
%    \begin{macrocode}
%<*ins>
\ifx\tudcdbatchfilemode\undefined%
%</ins>
%<*preamble>
\input{regression-test.tex}
%</preamble>
%<*ins>
\fi
%</ins>
%    \end{macrocode}
%
% \begin{unittest}{dev-version}{Kompilation der Artikelklasse}
%
% In dem ersten Test wird eine erfolgreiche Kompilation der Artikelklasse
% getestet.
%
%    \begin{macrocode}
%<*ins>
\generate{\file{\jobname-0001.lvt}{\from{\jobname.dtx}{0001,preamble}}}
%</ins>
%    \end{macrocode}
% Und dies ohne jegliche zusätzliche Pakete.
%    \begin{macrocode}
%<*ins>
\ifx\tudcdbatchfilemode\undefined%
%</ins>
%<*0001>
\documentclass{tudcdartcl}

\begin{document}
Hallo Welt!
\end{document}
%</0001>
%<*ins>
\fi
%</ins>
%    \end{macrocode}
%
% \end{unittest}
%
% \begin{unittest}{dev-version}{Kompilation der Reportklasse}
%
% In diesem Test wird eine erfolgreiche Kompilation der Reportklasse
% getestet.
%    \begin{macrocode}
%<*ins>
\generate{\file{\jobname-0002.lvt}{\from{\jobname.dtx}{0002,preamble}}}
%</ins>
%    \end{macrocode}
% Und dies ohne jegliche zusätzliche Pakete.
%    \begin{macrocode}
%<*ins>
\ifx\tudcdbatchfilemode\undefined%
%</ins>
%<*0002>
\documentclass{tudcdreprt}

\begin{document}
Hallo Welt!
\end{document}
%</0002>
%<*ins>
\fi
%</ins>
%    \end{macrocode}
% \end{unittest}
%
% \begin{unittest}{dev-version}{Kompilation der Buchklasse}
%
% In diesem Test wird eine erfolgreiche Kompilation der Buchklasse
% getestet.
%    \begin{macrocode}
%<*ins>
\generate{\file{\jobname-0003.lvt}{\from{\jobname.dtx}{0003,preamble}}}
%</ins>
%    \end{macrocode}
% Und dies ohne jegliche zusätzliche Pakete.
%    \begin{macrocode}
%<*ins>
\ifx\tudcdbatchfilemode\undefined%
%</ins>
%<*0003>
\documentclass{tudcdbook}

\begin{document}
Hallo Welt!
\end{document}
%</0003>
%<*ins>
\fi
%</ins>
%    \end{macrocode}
%
% \end{unittest}
%
% \begin{unittest}{dev-version}{Kompilation der Briefklasse}
%
% In diesem Test wird eine erfolgreiche Kompilation der Briefklasse
% getestet.
%    \begin{macrocode}
%<*ins>
\generate{\file{\jobname-0004.lvt}{\from{\jobname.dtx}{0004,preamble}}}
%</ins>
%    \end{macrocode}
% Und dies ohne jegliche zusätzliche Pakete.
%    \begin{macrocode}
%<*ins>
\ifx\tudcdbatchfilemode\undefined%
%</ins>
%<*0004>
\documentclass{tudcdlttr}

\begin{document}
Hallo Welt!
\end{document}
%</0004>
%<*ins>
\fi
%</ins>
%    \end{macrocode}
%
% \end{unittest}
%
% \begin{unittest}{dev-version}{Kompilation der Beamerklasse}
%
% In diesem Test wird eine erfolgreiche Kompilation der Beamerklasse
% getestet.
%    \begin{macrocode}
%<*ins>
\generate{\file{\jobname-0005.lvt}{\from{\jobname.dtx}{0005,preamble}}}
%</ins>
%    \end{macrocode}
% Und dies ohne jegliche zusätzliche Pakete.
%    \begin{macrocode}
%<*ins>
\ifx\tudcdbatchfilemode\undefined%
%</ins>
%<*0005>
\documentclass{beamer}

\usetheme{tudcd}

\begin{document}
\begin{frame}
  \frametitle{Beamertest}
  \framesubtitle{Beamertest subtitle}
  Hello World!
\end{frame}
\end{document}
%</0005>
%<*ins>
\fi
%</ins>
%    \end{macrocode}
%
% \end{unittest}
%
%    \begin{macrocode}
%<*ins>
\endbatchfile%
%</ins>
%    \end{macrocode}
%
